\section{Discussion}\label{sec:discussion}

The audit relocates the tail-small loss. Under standard overlapping tiling the slice boundary removes the pixels of essentially no small object, and for the audited tail-small set the loss occurs before any candidate is formed rather than in the merge, the score threshold or the output cap, so the admissible fix is a different detector or more training rather than a smarter read-out, with the important caveat that the dominant stage is contingent on the detector and the image source.

A plausible mechanism, which we state as inference rather than as a measured cause, concerns the training signal available to rare classes at small scales. A one-stage detector forms a candidate for an object only where it has learned a confident, correctly classified response; tail classes contribute few positive examples, and their small instances fewer still, so the detector may never fire a same-class box on a rare small pest in any view. The oracle is consistent with this reading: for head-small objects the full-view evidence is present and merely suppressed, so a better read-out recovers up to 8.3 points, whereas for tail-small objects the evidence is largely absent and the same read-out recovers a single object. The cross-detector contrast sharpens the mechanism rather than contradicting it. The undertrained FCOS detector actually leaves fewer tail-small objects without a candidate than YOLO11n, but only because it floods each tile with low-score boxes that the union merge then fuses away; the candidates it forms are not confident detections, and they are lost immediately downstream. Candidate formation in the sense that matters, a usable same-class response, remains the binding constraint for the tail.

These findings sit alongside, not against, the methods they test. Cross-view agreement rescoring genuinely improves tiled detection by re-ranking boxes that exist across tiles~\citep{na2026crossview}, and our partition explains precisely its reach: it recovers post-hoc-recoverable misses, which for head-small objects are plentiful but for tail-small objects number a handful. Calibration and prior-correction methods for long-tailed detection~\citep{pan2021norcal,alexandridis2025fracal} and output-cap analyses~\citep{dave2021devil} likewise operate on the decoded list and, as the oracle requires, cannot reach the no-candidate tail-small mass. Our ledger extends output-only error decomposition~\citep{hoiem2012diagnosing,bolya2020tide} one stage upstream, to candidate formation, which in this audit is where the tail-small outcome appears to be decided. That slicing on these downscaled images raises recall at the operating point without raising average precision is consistent with analyses of when sliced inference pays off as a function of resolution~\citep{becker2026scahi}, and the near-28-point leakage we document reinforces the warning that near-duplicate handling materially changes long-tailed benchmarks~\citep{nguyen2026ip102taildet}.

Four limitations bound what these results claim. The confirmatory endpoint is 38 tail-small objects in 27 groups, 25 of them a single leafhopper species, so the primary claim is a pooled three-category proportion and the per-class cells are descriptive only; we make no per-class powered statement, and a different tail composition could shift the pooled profile. Each detector is trained once, so we report the partition's stability across read-outs, tilings, thresholds and caps but not across training seeds, and the seed-to-seed variance of the 16-of-27 figure is unmeasured. RicePest-30 images are downscaled to a 640-pixel long edge, so ``small'' is defined in that frame and the candidate-formation finding is a statement about that regime; the native-resolution re-check is the deliberate counterweight, and it shows the loss moving to the read-out and the cap when the source changes. Finally, because the dominant stage depends jointly on the detector and the image source, the result is not a universal law of sliced inference but a diagnosis that must be re-run per deployment, which the frozen audit is designed to make straightforward. One scope limit deserves emphasis: the audit identifies the class of fix that is admissible for the tail-small no-candidate misses, but it does not demonstrate that any particular detector or training change recovers them, which would require the retraining this frozen-detector study deliberately excludes.

Within these bounds the implication for practice is concrete. For long-tailed small pests on downscaled imagery, the only admissible lever for the tail-small no-candidate misses is the detector and its training signal for rare small classes rather than the overlap, the merge rule or a post-hoc calibrator, so a team that invests in boundary-aware merging or score recalibration for this problem would be addressing only the minority of recoverable misses. For dense native-resolution traps the balance is different: the read-out and the output cap carry the loss, and replacing the union merge with NMS and lifting the cap recovers most of it. The general recommendation is therefore to run a one-pass recovery-ceiling audit before choosing a remedy, because it requires only a frozen inference pass and the bookkeeping ledger, and the partition it returns rules entire families of fixes in or out before any of them is built.

